%%%%%%%%%%%%%%%%%%%%%%%%%%%%%%%%%%%%%%%%%%%%%%%%%%%%%%%%%%%%%%%%%%%%%%
% Lớp 8 | Chương III. Tứ giác | Hình thoi và hình vuông
% Dạng 8H3?14-1: Tính chất, dấu hiệu nhận biết hình thoi
%%%%%%%%%%%%%%%%%%%%%%%%%%%%%%%%%%%%%%%%%%%%%%%%%%%%%%%%%%%%%%%%%%%%%%


% nguồn: De_AI_Toan8_cuoi_hoc_ki_i_1 (nhập 21/11/2026)
\begin{bt}%[8H3H14-1]%[ma:a6e92e83]
	Cho hình thoi $ABCD$ có độ dài hai đường chéo là $AC = 12\text{ cm}$ và $BD = 16\text{ cm}$. Gọi $O$ là giao điểm của hai đường chéo.
a) Tính độ dài cạnh của hình thoi và chu vi của hình thoi $ABCD$.
b) Tính diện tích của hình thoi $ABCD$.
	\loigiai{
		a) Trong hình thoi $ABCD$, hai đường chéo vuông góc với nhau tại trung điểm $O$ của mỗi đường, nên: \\
$OA = \dfrac{AC}{2} = 6\text{ cm}$, $OB = \dfrac{BD}{2} = 8\text{ cm}$ và $AC \perp BD$. \\
Áp dụng định lí Pythagore trong tam giác $OAB$ vuông tại $O$: $AB^2 = OA^2 + OB^2 = 6^2 + 8^2 = 100 \Rightarrow AB = 10\text{ cm}$. \\
Chu vi hình thoi $ABCD$ là $4 \cdot AB = 4 \cdot 10 = 40\text{ cm}$. \\
b) Diện tích hình thoi $ABCD$ là $S = \dfrac{1}{2} \cdot AC \cdot BD = \dfrac{1}{2} \cdot 12 \cdot 16 = 96\text{ cm}^2$.
	}
\end{bt}

% nguồn: De_AI_Toan8_cuoi_hoc_ki_i_1 (nhập 21/11/2026)
\begin{bt}%[8H3H14-1]%[ma:48e1029f]
	Cho hình thoi $ABCD$ có độ dài cạnh bằng $13\text{ cm}$ và đường chéo $AC = 10\text{ cm}$.
a) Tính độ dài đường chéo $BD$.
b) Tính diện tích của hình thoi $ABCD$.
	\loigiai{
		a) Gọi $O$ là giao điểm hai đường chéo $AC$ và $BD$. Ta có $AC \perp BD$ tại $O$ và $OA = \dfrac{AC}{2} = 5\text{ cm}$. \\
Áp dụng định lí Pythagore trong tam giác $OAB$ vuông tại $O$: $OB^2 = AB^2 - OA^2 = 13^2 - 5^2 = 169 - 25 = 144 \Rightarrow OB = 12\text{ cm}$. \\
Suy ra $BD = 2OB = 2 \cdot 12 = 24\text{ cm}$. \\
b) Diện tích hình thoi $ABCD$ là $S = \dfrac{1}{2} \cdot AC \cdot BD = \dfrac{1}{2} \cdot 10 \cdot 24 = 120\text{ cm}^2$.
	}
\end{bt}

% nguồn: De_AI_Toan8_cuoi_hoc_ki_i_1 (nhập 21/11/2026)
\begin{bt}%[8H3H14-1]%[ma:52cd0e33]
	Cho hình thoi $ABCD$ có cạnh $AB = 8\text{ cm}$ và $\widehat{A} = 60^\circ$.
a) Chứng minh tam giác $ABD$ là tam giác đều, từ đó tính độ dài đường chéo $BD$.
b) Tính độ dài đường chéo $AC$.
	\loigiai{
		a) Vì $ABCD$ là hình thoi nên $AB = AD = 8\text{ cm}$, suy ra tam giác $ABD$ cân tại $A$. \\
Lại có $\widehat{A} = 60^\circ$ nên tam giác $ABD$ là tam giác đều, do đó $BD = AB = 8\text{ cm}$. \\
b) Gọi $O$ là giao điểm hai đường chéo $AC$ và $BD$. Ta có $OB = \dfrac{BD}{2} = 4\text{ cm}$ và $AC \perp BD$. \\
Áp dụng định lí Pythagore trong tam giác $OAB$ vuông tại $O$: $OA^2 = AB^2 - OB^2 = 8^2 - 4^2 = 48 \Rightarrow OA = 4\sqrt{3}\text{ cm}$. \\
Vậy $AC = 2OA = 8\sqrt{3}\text{ cm}$.
	}
\end{bt}

% nguồn: De_AI_Toan8_cuoi_hoc_ki_i_1 (nhập 21/11/2026)
\begin{bt}%[8H3H14-1]%[ma:1fa82fea]
	Cho tam giác $ABC$ cân tại $A$, gọi $M$ là trung điểm của cạnh $BC$. Lấy điểm $D$ đối xứng với $A$ qua $M$.
a) Chứng minh tứ giác $ABDC$ là hình thoi.
b) Chứng minh $AD$ là tia phân giác của góc $BAC$.
	\loigiai{
		a) Tứ giác $ABDC$ có hai đường chéo $AD$ và $BC$ cắt nhau tại trung điểm $M$ của mỗi đường nên $ABDC$ là hình bình hành. \\
Lại có $AB = AC$ (vì tam giác $ABC$ cân tại $A$), nên hình bình hành $ABDC$ có hai cạnh kề bằng nhau, do đó $ABDC$ là hình thoi. \\
b) Vì $ABDC$ là hình thoi nên đường chéo $AD$ là đường phân giác của góc $BAC$.
	}
\end{bt}

% nguồn: De_AI_Toan8_cuoi_hoc_ki_i_1 (nhập 21/11/2026)
\begin{bt}%[8H3H14-1]%[ma:ee6abab3]
	Cho hình chữ nhật $ABCD$ có $AB = 8\text{ cm}$, $BC = 6\text{ cm}$. Gọi $M$, $N$, $P$, $Q$ lần lượt là trung điểm của các cạnh $AB$, $BC$, $CD$, $DA$.
a) Chứng minh tứ giác $MNPQ$ là hình thoi.
b) Tính độ dài cạnh $MN$ và diện tích của hình thoi $MNPQ$.
	\loigiai{
		a) Tương tự bài toán đường trung bình trong các tam giác $ABC$, $ADC$, $ABD$, $CBD$, ta có $MN = PQ = \dfrac{1}{2}AC$ và $MQ = NP = \dfrac{1}{2}BD$. \\
Vì $ABCD$ là hình chữ nhật nên $AC = BD$, suy ra $MN = NP = PQ = QM$. Vậy tứ giác $MNPQ$ là hình thoi. \\
b) Xét tam giác $MBN$ vuông tại $B$ có $MB = \dfrac{AB}{2} = 4\text{ cm}$ và $BN = \dfrac{BC}{2} = 3\text{ cm}$. \\
Theo định lí Pythagore: $MN^2 = MB^2 + BN^2 = 4^2 + 3^2 = 25 \Rightarrow MN = 5\text{ cm}$. \\
Hai đường chéo của hình thoi $MNPQ$ là $MP = BC = 6\text{ cm}$ và $NQ = AB = 8\text{ cm}$, nên diện tích $S_{MNPQ} = \dfrac{1}{2} \cdot 6 \cdot 8 = 24\text{ cm}^2$.
	}
\end{bt}

% nguồn: chuyen_dang (nhập 22/11/2026)
\begin{bt}%[8H3H14-1]%[ma:66b19c4f]
	Cho hình thoi $ABCD$ cạnh $2$ và $\widehat{BAD}=60^\circ$.  Tìm độ dài $\overrightarrow{AC}$.
	\loigiai{\begin{center}
		\begin{tikzpicture}[scale=1,line join=round,line cap=round,font=\footnotesize,>=stealth]
				 \def\a{3}
				\def\b{5.25} 
				\path  (0:0) coordinate (B)
					++(0:\a) coordinate (D) 
					([turn]120:\a)  coordinate (A)
					([turn]-150:\b) coordinate  (C)
					($(B)!0.5!(D)$)  coordinate (O);
				\draw (D)--(A)--(B)--(D)--(C)--(B) 
					(A)--(C); 
				\foreach \x/ \goc  in {A/135,B/-135,D/-45,C/-30,O/40}
				\fill (\x) circle (1pt) 
					 ($(\x)+(\goc:3mm)$) node {$\x$};
			\end{tikzpicture}
	\end{center}
	 Gọi $O$ là giao điểm của hai đường chéo $AC$ và $BD$.\\
	Tam giác $ABD$ có $AB=AD$ và $\widehat{BAD}=60^\circ$ nên $\Delta ABD$ đều canh bằng $2$\\ 
	$\Rightarrow  AO=\dfrac{2\sqrt{3}}{2}=\sqrt{3}$.\\
	Khi đó $\left| \overrightarrow{AC}  \right| =AC=2AO=2\sqrt3 \approx 3{,}46$.
 }
\end{bt}

% nguồn: chuyen_dang (nhập 22/11/2026)
\begin{bt}%[8H3H14-1]%[ma:66e4858f]
	Hình bình hành có hai cạnh là $5$ và $9$, một đường chéo bằng $11$.  Tìm độ dài đường chéo còn lại
	\loigiai{ 
		\begin{center}
			\begin{tikzpicture}[>=stealth,line join=round,line cap=round]
					\draw  (0,0) coordinate (A)--(0:4) coordinate (B)--($(80:2)+(0:4)$) coordinate (C)--(80:2) coordinate (D)--cycle;
					\foreach \diem/\vitrin in {A/left,B/right,C/right,D/left} \fill[black](\diem)  circle (1.5pt) node[\vitrin]{$\diem$};
					\path (A)--(B)  node[pos=0.5,sloped,below]{$9$};
					\path  (A)--(D) node[pos=0.5,left]{$5$};
					 \path (D)--(B) node[pos=0.55,sloped,below]{11};
					\draw  (A)--(C)
						 (B)--(D);
				\end{tikzpicture}
		\end{center}
		Gọi hình bình  hành là $ABCD$, $AD=5$, $AB=9$.\\
		Gọi  $\alpha$ là góc đối diện với đường chéo có độ dài $11$.\\
		Ta có $\cos \alpha=\dfrac{5^2+9^2-11^2}{2\cdot  5\cdot 9}=-\dfrac{1}{6}$
		$\Rightarrow \alpha$ là góc tù $\Rightarrow \alpha=\widehat{BAD}$ $\Rightarrow  BD=11$
		 \[\Rightarrow AC^2=AD^2+DC^2-2\cdot AD\cdot DC\cdot \cos \widehat{ADC}=AD^2+DC^2+2\cdot AD\cdot DC\cdot \cos \widehat{BAD}\]
		$\widehat{BAD}$ và $\widehat{ADC}$  bù nhau $\Rightarrow \cos \widehat{ADC}=-\cos \widehat{BAD}$
		\[\Rightarrow AC^2=5^2+9^2+2\cdot 5\cdot  9\cdot \left(-\dfrac{1}{6} \right)=91\Rightarrow AC=\sqrt{91}.\]
	} 
\end{bt}
